% !TeX root = 04_intro_ml_polytech.tex
% !TeX encoding = UTF-8
% !TeX spellcheck = en_GB
\begin{frame}{Practical Neural Networks}{Some peculiarities of deep learning\dots}
    Neural network training often uses \textbf{GPU}s which are faster than regular CPUs:

    \vfill

    \begin{itemize}
        \item GPUs have \emph{usually} \textbf{less memory} than CPUs
        \item Data sets can contain millions of large files (e.g.\ videos, images, text)
        \item Data cannot be fed directly into GPU memory but needs to be split into \textbf{batches}
    \end{itemize}

    \vfill

    \begin{center}
        \begin{tikzpicture}[
        block/.style={thick, CEAred, fill=CEAred!15},
        block blue/.style={thick, CEAblue, fill=CEAblue!15},
        output/.style={circle, draw=CEAgreen, fill=CEAgreen!15},
        direct/.style={ultra thick, CEAblack, ->}
    ]

    % Data
    \draw<1>[block] (0, 0) rectangle (2, 3) node[midway] {DATA};
    \draw<2->[block, opacity=0.5] (0, 0) rectangle (2, 0.4) node[midway] {BATCH$_1$};
    \draw<2->[block, opacity=0.5] (0, 0.5) rectangle (2, 0.9) node[midway] {BATCH$_2$};
    \draw<2->[block] (0, 1.0) rectangle (2, 1.4) node[midway] {BATCH$_3$};
    \draw<2->[block, opacity=0.5] (0, 1.5) rectangle (2, 1.9) node[midway] {BATCH$_4$};
    \draw<2->[block, opacity=0.5] (0, 2.0) rectangle (2, 2.4) node[midway] {BATCH$_5$};
    \draw<2->[block, opacity=0.5] (0, 2.5) rectangle (2, 3.0) node[midway] {BATCH$_6$};

    % Model
    \draw<1->[block blue] (4.0, 1.0) rectangle (6, 2.0) node[midway] {MODEL};
    \node<1->[output] (o) at (8.3, 1.5) {$\hat{y}$};

    % Arrows
    \draw<1>[direct] (2.1, 1.5) -- (3.9, 1.5) node[above, midway, sloped] {feed};
    \draw<2->[direct] (2.1, 1.2) -- (3.9, 1.5) node[above, midway, sloped] {feed};

    \draw<1->[direct] (6.1, 1.5) -- (7.9, 1.5) node[above, midway, sloped] {compute};
    \draw<1->[direct] (o.south) to[bend left] (5.0, 0.9);
    \node<1>[anchor=north] at (6.5, 0.5) {gradient descent};
    \node<2->[anchor=north] at (6.5, 0.5) {\alert{mini-batch} gradient descent};
\end{tikzpicture}

    \end{center}

    \vfill
    {\tiny You can technically have a number of batches equal to the number of data points, in which case you talk about \textbf{stochastic gradient descent}.}
\end{frame}

\begin{frame}[fragile]{Practical Neural Networks}{PyTorch simplifies all\dots}
    The library we chose (\textbf{PyTorch}) already contains some helpers:

    \vfill

    \begin{onlyenv}<1>
        Create a \textbf{dataset} from vectors/tensors:
        \begin{codewithcard}{Create a dataset}
import torch

X = np.array(...)
y = np.array(...)

# Convert to PyTorch objects
X = torch.from_numpy(X)
y = torch.from_numpy(y)

# Create a dataset that can be fed to a PyTorch model
dataset = torch.utils.data.TensorDataset(X, y)
        \end{codewithcard}
    \end{onlyenv}

    \begin{onlyenv}<2>
        Split a dataset into a generator of \textbf{batches}:
        \begin{codewithcard}{Split into batches}
batches = torch.utils.data.DataLoader(dataset,
                                      batch_size=32,
                                      shuffle=True,  # train: T, val/test: F
                                      drop_last=True,  # train: T, val/test: F
                                    )
        \end{codewithcard}

        \vfill
        {\tiny It is customary to use powers of 2 for the batch size.}
    \end{onlyenv}

    \begin{onlyenv}<3->
        Automatically choose the device (CPU or GPU) to execute the computations:
        \begin{codewithcard}{Select the device}
device = torch.device("cuda" if torch.cuda.is_available() else "cpu")
        \end{codewithcard}
    \end{onlyenv}
\end{frame}

\begin{frame}[fragile]{Practical Neural Networks}{PyTorch simplifies all\dots Almost\dots}
    You are then left with the task of handling the model and training \textbf{loop}:

    \vfill

    \begin{onlyenv}<1>
        \begin{codewithcard}{Model definition}
in_dimension = ...; out_dimension = ...

model = torch.nn.Sequential(
            torch.nn.Linear(in_dimension, 8),
            torch.nn.ReLU(),
            torch.nn.Linear(8, 4),
            torch.nn.ReLU(),
            torch.nn.Linear(4, out_dimension),
        ).to(device)
        \end{codewithcard}

        \vfill

        {\tiny Different types of layers and activations are available in the \texttt{torch.nn} module. \texttt{Linear} layers take tabular data of size $N \times P$ directly (convolutional layers need an explicit channel axis, see later).}
    \end{onlyenv}

    \begin{onlyenv}<2>
        \begin{codewithcard}{Loss and optimiser}
criterion = torch.nn.MSELoss()
optim = torch.optim.AdamW(model.parameters(), lr=0.001)
        \end{codewithcard}

        \vfill

        {\tiny \texttt{AdamW} is a powerful gradient descent algorithm, widely used in deep learning. You can however experiment with other optimisers.}
    \end{onlyenv}


    \begin{onlyenv}<3>
        \begin{codewithcard}{Training loop}
for n in range(epochs):  # see the dataset multiple times
    for X, y in batches:  # batches created before
            # Move to device
            X = X.to(device); y = y.to(device)

            # Compute predictions
            optim.zero_grad(); model.train(True)
            y_pred = model(X)

            # Compute the loss and optimiser step
            loss = criterion(y_pred, y); loss.backward()
            optim.step()
        \end{codewithcard}
    \end{onlyenv}
\end{frame}

\begin{frame}{Practical Neural Networks}{PyTorch simplifies all\dots}
    Neural networks are very \textbf{versatile} and allow a lot more operations on data.
    Here is some important information to keep in mind:

    \vfill

    \begin{itemize}
        \item \textbf{Linear layers} (\texttt{torch.nn.Linear}): they are the classical \emph{fully-connected} layers (\textbf{MLP} layers)
        \item \textbf{Convolutional layers} (\texttt{torch.nn.ConvXd}): they compute a sliding \emph{convolution} of the input in $X$ dimensions (1, 2, or 3). Must specify kernel size
        \item \textbf{Activation layers} (e.g.\ \texttt{torch.nn.ReLU}): they are activation functions that need to be inserted between hidden layers
        \item \textbf{Dropout layers} (\texttt{torch.nn.Dropout}): \emph{dropout} is an operation to switch off some neurons during training for regularisation
        \item \textbf{Pooling layers} (e.g.\ \texttt{torch.nn.AdaptiveAvgPoolXd}): they are responsable of operations such as mean, max, etc. in $X$ dimensions (1, 2, or 3)
        \item \textbf{Flatten layer} (\texttt{torch.nn.Flatten}): useful to flatten some axis to make them compatible with other layers
    \end{itemize}
\end{frame}

\begin{frame}[fragile]{Hands-on Session Neural Networks}{\includegraphics[height=12pt]{01_slides/img/tool.pdf}~Building neural networks (I)\dots}
    \begin{itemize}
        \item \alert<2>{\includegraphics[height=12pt]{01_slides/img/pen-tool.pdf}~Open the \texttt{machine\_learning.ipynb} notebook and run all cells}
        \item \alert<3>{\includegraphics[height=12pt]{01_slides/img/pen-tool.pdf}~Determine what device (GPU or CPU) you are going to use}
        \item \alert<4>{\includegraphics[height=12pt]{01_slides/img/pen-tool.pdf}~Create a new Python file called \texttt{tools.py} inside the \texttt{src/mlintro} directory}
        \item \alert<5>{\includegraphics[height=12pt]{01_slides/img/pen-tool.pdf}~Define a new function \texttt{create\_dataset} in \texttt{tools.py} which takes features and labels (train, validation, or test) and returns the \texttt{DataLoader} of the batches}
        \item \alert<6>{\includegraphics[height=12pt]{01_slides/img/pen-tool.pdf}~Create the training, validation, and test batches in the notebook using the \textbf{PCA}-reduced data}
        \item \alert<7>{\includegraphics[height=12pt]{01_slides/img/pen-tool.pdf}~Define a \textbf{fully-connected} model to process the \textbf{PCA}-reduced data}
        \item \alert<8>{\includegraphics[height=12pt]{01_slides/img/pen-tool.pdf}~Choose an appropriate classification loss and an optimiser}
    \end{itemize}

    \begin{onlyenv}<3>
        \begin{tikzpicture}[overlay]
            \node[set, xshift=-2em, yshift=-5em, text width=35em] at (current page.center) {
                \begin{notebookin}
import torch
device = torch.device("cuda" if torch.cuda.is_available() else "cpu")
print("Device:", device)
                \end{notebookin}
                \begin{notebookout}
Device: cpu
                \end{notebookout}
            };
        \end{tikzpicture}
    \end{onlyenv}

    \begin{onlyenv}<5>
        \begin{tikzpicture}[overlay]
            \node[xshift=-2em, yshift=-6em, text width=35em] at (current page.center) {
                \begin{codewithcard}{tools.py}
import torch
def create_dataset(features, labels, batch_size=32, train=False):
    # Convert to tensors
    features = torch.from_numpy(features)
    labels = torch.from_numpy(labels)
    # Create a tensor dataset
    dataset = torch.utils.data.TensorDataset(features, labels)
    # Return the batches
    return torch.utils.data.DataLoader(dataset,
                                       batch_size=batch_size,
                                       shuffle=train,
                                       drop_last=train
                                      )
                \end{codewithcard}
            };
        \end{tikzpicture}
    \end{onlyenv}

    \begin{onlyenv}<6>
        \begin{tikzpicture}[overlay]
            \node[set, xshift=-2em, yshift=-6em, text width=35em] at (current page.center) {
                \begin{notebookin}
from mlintro.tools import create_dataset

# Use y-1 to have classes start from 0
train = create_dataset(X_train_pca, y_train_sel-1, train=True)
val = create_dataset(X_val_pca, y_val_sel-1, train=False)
test = create_dataset(X_test_pca, y_test_sel-1, train=False)
                \end{notebookin}
            };
        \end{tikzpicture}
    \end{onlyenv}

    \begin{onlyenv}<7>
        \begin{tikzpicture}[overlay]
            \node[set, xshift=-2em, yshift=-6em, text width=35em] at (current page.center) {
                \begin{notebookin}
in_dim = X_train_pca.shape[1]
out_dim = int(y_train_sel.max())
model = torch.nn.Sequential(
            torch.nn.Linear(in_dim, 4),
            torch.nn.ReLU(),
            torch.nn.Linear(4, 2),
            torch.nn.ReLU(),
            torch.nn.Linear(2, out_dim)
        ).to(device)
                \end{notebookin}
            };
        \end{tikzpicture}
    \end{onlyenv}

    \begin{onlyenv}<8>
        \begin{tikzpicture}[overlay]
            \node[set, xshift=-2em, yshift=-6em, text width=35em] at (current page.center) {
                \begin{notebookin}
criterion = torch.nn.CrossEntropyLoss()
optim = torch.optim.AdamW(model.parameters(), lr=0.001)
                \end{notebookin}
            };
        \end{tikzpicture}
    \end{onlyenv}
\end{frame}

\begin{frame}[fragile]{Hands-on Session Neural Networks}{\includegraphics[height=12pt]{01_slides/img/tool.pdf}~Building neural networks (II)\dots}
    \begin{itemize}
        \item \alert<2>{\includegraphics[height=12pt]{01_slides/img/pen-tool.pdf}~Define a function in \texttt{tools.py} that \textbf{trains} a single batch $(X, y)$ of data (the function must return the value of the loss and the accuracy of the batch)}
        \item \alert<3>{\includegraphics[height=12pt]{01_slides/img/pen-tool.pdf}~Define a function in \texttt{tools.py} that uses the previous to train a model for an entire \textbf{epoch}}
        \item \alert<4-5>{\includegraphics[height=12pt]{01_slides/img/pen-tool.pdf}~Define the previous functions in \texttt{tools.py} but for the \textbf{validation set}}
        \item \alert<6>{\includegraphics[height=12pt]{01_slides/img/pen-tool.pdf}~Define a function in \texttt{tools.py} called \texttt{trainer} which \textbf{trains and validates} the model for a given number of epochs. The function must return the list of train/val losses and accuracies per epoch}
        \item \alert<7>{\includegraphics[height=12pt]{01_slides/img/pen-tool.pdf}~Define a function in \texttt{tools.py} which \textbf{predicts} the classes (without computing the loss) for a single batch, then a function which \emph{concatenates} the prediction for all batches in the data loader}
    \end{itemize}

    \begin{onlyenv}<2>
        \begin{tikzpicture}[overlay]
            \node[xshift=-2em, yshift=-6em, text width=35em] at (current page.center) {
                \begin{codewithcard}{tools.py}
from sklearn.metrics import accuracy_score
def train_single_batch(X, y, model, optim, criterion, device):
    # Put on device
    X = X.to(device); y = y.to(device)
    # Compute predictions
    optim.zero_grad(); model.train(True)
    y_pred = model(X)
    # Compute loss
    loss = criterion(y_pred, y); loss.backward(); optim.step()
    # Compute accuracy
    acc = float(accuracy_score(y.cpu(), y_pred.argmax(dim=1).cpu()))
    return loss.item(), acc
                \end{codewithcard}
            };
        \end{tikzpicture}
    \end{onlyenv}

    \begin{onlyenv}<3>
        \begin{tikzpicture}[overlay]
            \node[xshift=-3em, yshift=-6em, text width=40em] at (current page.center) {
                \begin{codewithcard}{tools.py}
def train_single_epoch(batches, model, optim, criterion, device):
    # Unpack the batches
    epoch_loss = 0.0; epoch_acc = 0.0
    for X, y in batches:
        batch_loss, batch_acc = train_single_batch(X, y, model, optim, criterion, device)
        epoch_loss += batch_loss; epoch_acc += batch_acc
    return epoch_loss / len(batches), epoch_acc / len(batches)
                \end{codewithcard}
            };
        \end{tikzpicture}
    \end{onlyenv}

    \begin{onlyenv}<4>
        \begin{tikzpicture}[overlay]
            \node[xshift=-2em, yshift=-6em, text width=35em] at (current page.center) {
                \begin{codewithcard}{tools.py}
def val_single_batch(X, y, model, criterion, device):
    # Put on device
    X = X.to(device); y = y.to(device)
    # Compute predictions
    model.train(False)
    with torch.no_grad():  # no gradients needed
        y_pred = model(X)
    # Compute loss
    loss = criterion(y_pred, y)
    # Compute accuracy
    acc = float(accuracy_score(y.cpu(), y_pred.argmax(dim=1).cpu()))
    return loss.item(), acc
                \end{codewithcard}
            };
        \end{tikzpicture}
    \end{onlyenv}

    \begin{onlyenv}<5>
        \begin{tikzpicture}[overlay]
            \node[xshift=-3em, yshift=-6em, text width=40em] at (current page.center) {
                \begin{codewithcard}{tools.py}
def val_single_epoch(batches, model, criterion, device):
    # Unpack the batches
    epoch_loss = 0.0; epoch_acc = 0.0
    for X, y in batches:
        batch_loss, batch_acc = val_single_batch(X, y, model, criterion, device)
        epoch_loss += batch_loss; epoch_acc += batch_acc
    return epoch_loss / len(batches), epoch_acc / len(batches)
                \end{codewithcard}
            };
        \end{tikzpicture}
    \end{onlyenv}

    \begin{onlyenv}<6>
        \begin{tikzpicture}[overlay]
            \node[xshift=-3em, yshift=-6em, text width=42em] at (current page.center) {
                \begin{codewithcard}{tools.py}
def trainer(n_epochs, train_batches, val_batches, model, optim, criterion, device):
    train_losses, train_accs, val_losses, val_accs = [], [], [], []
    for n in range(n_epochs):
        train_loss, train_acc = train_single_epoch(train_batches, model, optim, criterion, device)
        train_losses.append(train_loss); train_accs.append(train_acc)
        val_loss, val_acc = val_single_epoch(val_batches, model, criterion, device)
        val_losses.append(val_loss); val_accs.append(val_acc)
    return train_losses, val_losses, train_accs, val_accs
                \end{codewithcard}
            };
        \end{tikzpicture}
    \end{onlyenv}

    \begin{onlyenv}<7>
        \begin{tikzpicture}[overlay]
            \node[xshift=-3em, yshift=-6em, text width=40em] at (current page.center) {
                \begin{codewithcard}{tools.py}
def predict_single_batch(X, model, device):
    X = X.to(device); model.train(False)
    with torch.no_grad():
        y_pred = model(X)
    return y_pred
def predict(batches, model, device):
    pred_list = []
    for X, _ in batches:
        pred = predict_single_batch(X, model, device)
        pred_list.append(pred)
    return torch.cat(pred_list)
                \end{codewithcard}
            };
        \end{tikzpicture}
    \end{onlyenv}
\end{frame}

\begin{frame}[fragile]{Hands-on Session Neural Networks}{\includegraphics[height=12pt]{01_slides/img/tool.pdf}~Building neural networks (III)\dots}
    \begin{itemize}
        \item \alert<2>{\includegraphics[height=12pt]{01_slides/img/pen-tool.pdf}~Train the fully-connected model using the \texttt{trainer} function}
        \item \alert<3>{\includegraphics[height=12pt]{01_slides/img/pen-tool.pdf}~Plot the losses and accuracy as functions of the number of epochs using \texttt{matplotlib}}
        \item \alert<4>{\includegraphics[height=12pt]{01_slides/img/pen-tool.pdf}~Use the prediction function to predict the classes of the test set, then compute the accuracy on the test set}
    \end{itemize}

    \begin{onlyenv}<2>
        \begin{tikzpicture}[overlay]
            \node[set, xshift=-2em, yshift=-14em, text width=35em] at (current page.center) {
                \begin{notebookin}
from mlintro.tools import trainer

(train_losses,
 val_losses,
 train_accs,
 val_accs) = trainer(25, train, val, model, optim, criterion, device)
                \end{notebookin}
            };
        \end{tikzpicture}
    \end{onlyenv}

    \begin{onlyenv}<3>
        \begin{tikzpicture}[overlay]
            \node[set, xshift=-2em, yshift=-13em, text width=35em] at (current page.center) {
                \centering
                \includegraphics[height=0.5\textheight]{04_slides/img/training_curves.png}
                \hfil
                \includegraphics[height=0.5\textheight]{04_slides/img/training_curves_acc.png}
            };
        \end{tikzpicture}
    \end{onlyenv}

    \begin{onlyenv}<4>
        \begin{tikzpicture}[overlay]
            \node[set, xshift=-2em, yshift=-13em, text width=35em] at (current page.center) {
                \begin{notebookin}
logits_pred = predict(test, model, device)
y_pred = logits_pred.argmax(dim=1).cpu()

acc = accuracy_score(y_test_sel-1, y_pred)
print(f"Test accuracy: {acc:.1%}")
                \end{notebookin}
                \begin{notebookout}
Test accuracy: 71.8%
                \end{notebookout}
            };
        \end{tikzpicture}
    \end{onlyenv}
\end{frame}

\begin{frame}[fragile]{Hands-on Session Neural Networks}{\includegraphics[height=12pt]{01_slides/img/tool.pdf}~Building neural networks (IV - optional)\dots}
    \begin{itemize}
        \item \alert<2>{\includegraphics[height=12pt]{01_slides/img/pen-tool.pdf}~Reshape the original data (the \enquote{selected} data, not the PCA-reduced ones) to have $1$ channel (i.e.\ one additional dimension as second axis: shape from $N \times P$ to $N \times 1 \times P$), then create \textbf{new} batches from them}
        \item \alert<3>{\includegraphics[height=12pt]{01_slides/img/pen-tool.pdf}~Use the \texttt{torch.nn.Conv1d} layer to define a convolutional model. Combine the \texttt{torch.nn.AdaptiveAvgPool1d} layer to compute the mean of the spectral dimension and the \texttt{torch.nn.Flatten} layer to make the output shape compatible with the previous model}
        \item \alert<4>{\includegraphics[height=12pt]{01_slides/img/pen-tool.pdf}~As before, train the model and display the loss and accuracy}
        \item \alert<5>{\includegraphics[height=12pt]{01_slides/img/pen-tool.pdf}~As before, predict the classes on the test set}
        \item \alert<6>{\includegraphics[height=12pt]{01_slides/img/brain.pdf}~Which model performs best for you?}
    \end{itemize}

    \begin{onlyenv}<2>
        \begin{tikzpicture}[overlay]
            \node[set, xshift=-2em, yshift=-10em, text width=38em] at (current page.center) {
                \begin{notebookin}
P = X_train_sel.shape[1]
X_train_sel_conv = X_train_sel.reshape(-1, 1, P)  # 1 channel input
X_val_sel_conv = X_val_sel.reshape(-1, 1, P)  # 1 channel input
X_test_sel_conv = X_test_sel.reshape(-1, 1, P)  # 1 channel input
train_conv = create_dataset(X_train_sel_conv, y_train_sel-1, train=True)
val_conv = create_dataset(X_val_sel_conv, y_val_sel-1, train=False)
test_conv = create_dataset(X_test_sel_conv, y_test_sel-1, train=False)
                \end{notebookin}
            };
        \end{tikzpicture}
    \end{onlyenv}

    \begin{onlyenv}<3>
        \begin{tikzpicture}[overlay]
            \node[set, xshift=-2em, yshift=-8em, text width=35em] at (current page.center) {
                \begin{notebookin}
model = torch.nn.Sequential(
            torch.nn.Conv1d(1, 8, kernel_size=5, stride=5),
            torch.nn.ReLU(),
            torch.nn.Conv1d(8, 16, kernel_size=3, stride=3),
            torch.nn.ReLU(),
            torch.nn.Conv1d(16, out_dim, kernel_size=3),  # logits
            torch.nn.AdaptiveAvgPool1d(1),
            torch.nn.Flatten(1)  # must be compatible with N x P
        ).to(device)
# new model -> new optimiser
optim = torch.optim.AdamW(model.parameters(), lr=0.001)
                \end{notebookin}
            };
        \end{tikzpicture}
    \end{onlyenv}

    \begin{onlyenv}<4>
        \begin{tikzpicture}[overlay]
            \node[set, xshift=-2em, yshift=-7em, text width=35em] at (current page.center) {
                \centering
                \includegraphics[height=0.5\textheight]{04_slides/img/training_curves_conv.png}
                \hfil
                \includegraphics[height=0.5\textheight]{04_slides/img/training_curves_conv_acc.png}
            };
        \end{tikzpicture}
    \end{onlyenv}

    \begin{onlyenv}<5>
        \begin{tikzpicture}[overlay]
            \node[set, xshift=-2em, yshift=-8em, text width=35em] at (current page.center) {
                \begin{notebookin}
logits_pred = predict(test_conv, model, device)
y_pred = logits_pred.argmax(dim=1).cpu()

acc = accuracy_score(y_test_sel-1, y_pred)
print(f"Test accuracy: {acc:.1%}")
                \end{notebookin}
                \begin{notebookout}
Test accuracy: 80.1%
                \end{notebookout}
            };
        \end{tikzpicture}
    \end{onlyenv}
\end{frame}

\begin{frame}{Neural Networks}{Preparing for exam!}
    You should now be proficient in:
    \begin{itemize}
        \item What neural networks are and how they are trained
        \item Using simple neural network models to train a classifier
    \end{itemize}
    \vfill \pause
    \begin{center}
        \huge
        Finals are next!
        \\[1em]
        Are there questions?
    \end{center}
    \vfill
\end{frame}
